\section{LCA $O(1)$ con Euler tour y RMQ}
\whenuse{Árbol estático con muchísimas consultas de LCA. Usar binary lifting si
también se necesita subir $k$ niveles o ahorrar memoria.}

El Euler tour guarda un vértice al visitarlo y cada vez que el DFS regresa a él.
El LCA de $a$ y $b$ es el vértice de menor profundidad entre sus primeras
apariciones en el recorrido. Una Sparse Table responde este mínimo en $O(1)$.

\complexity{Preproceso $O(n\log n)$; consulta $O(1)$; memoria $O(n\log n)$.}

\begin{lstlisting}[style=cpp]
struct EulerRMQLCA {
    vector<int> first, depth, euler, logarithm;
    vector<vector<int>> table;
    EulerRMQLCA(const vector<vector<int>> &tree, int root = 0)
        : first(tree.size(), -1), depth(tree.size()) {
        assert(!tree.empty());
        dfs(root, -1, tree);
        int length = euler.size();
        logarithm.assign(length + 1, 0);
        for (int i = 2; i <= length; ++i)
            logarithm[i] = logarithm[i / 2] + 1;
        // table[k][i]: vertice de menor profundidad
        // en euler[i .. i + 2^k).
        table.push_back(euler);

        for (int k = 1; (1 << k) <= length; ++k) {
            int half = 1 << (k - 1);
            int count = length - (1 << k) + 1;
            table.push_back(vector<int>(count));
            for (int i = 0; i < count; ++i) {
                int a = table[k - 1][i];
                int b = table[k - 1][i + half];
                table[k][i] = depth[a] < depth[b] ? a : b;
            }
        }
    }
    void dfs(int v, int parent,
             const vector<vector<int>> &tree) {
        first[v] = euler.size();
        euler.push_back(v);
        for (int to : tree[v]) {
            if (to == parent) continue;
            depth[to] = depth[v] + 1;
            dfs(to, v, tree);
            euler.push_back(v); // Regreso desde el hijo.
        }
    }
    int lca(int a, int b) const {
        int left = first[a], right = first[b];
        if (left > right) swap(left, right);
        int k = logarithm[right - left + 1];
        int x = table[k][left];
        int y = table[k][right - (1 << k) + 1];

        return depth[x] < depth[y] ? x : y;
    }
    int distance(int a, int b) const {
        int ancestor = lca(a, b);
        return depth[a] + depth[b] - 2 * depth[ancestor];
    }
};

EulerRMQLCA lca(tree, 0);
cout << lca.lca(3, 4) << '\n';      // Ancestro comun
cout << lca.distance(3, 2) << '\n'; // Distancia en aristas
\end{lstlisting}


\warning{El grafo debe ser un árbol conectado. El DFS recursivo puede desbordar
la pila en árboles muy profundos.}